\chapter{Cel i zakres projektu}

\section{Wstęp}
Projekt koncentruje się na zaprojektowaniu oraz implementacji nowoczesnych interfejsów użytkownika dla wewnętrznego systemu zarządzania hotelem \textbf{French Riviera}. Celem opracowania jest stworzenie cyfrowego ekosystemu, który optymalizuje komunikację wewnątrz obiektu, automatyzuje obieg zgłoszeń oraz dostarcza kadrze zarządzającej precyzyjnych danych operacyjnych.

System został zaprojektowany w architekturze wielomodułowej, dostosowanej do zróżnicowanych urządzeń końcowych: od stacjonarnych paneli administracyjnych, przez tablety pokojowe i kuchenne, aż po natywne aplikacje mobilne dla personelu terenowego.

\section{Zakres projektu}
System \textbf{French Riviera} jest rozwiązaniem klasy \textit{Internal Management System}, co oznacza, że jego zakres kończy się na wewnętrznej infrastrukturze hotelu. Projekt pomija etap zewnętrznej rezerwacji miejsc (\textit{booking}) oraz publiczne systemy płatności, skupiając się w pełni na obsłudze gościa już zameldowanego oraz na efektywnym zarządzaniu zasobami ludzkimi i technicznymi obiektu.

\section{Cel projektu}
Głównym założeniem jest stworzenie spójnego środowiska GUI, które:
\begin{itemize}
    \item zapewnia gościom bezobsługowy dostęp do pełnego wachlarza usług hotelowych,
    \item rozdziela zadania operacyjne pomiędzy wyspecjalizowane działy (serwis, utrzymanie ruchu, gastronomia),
    \item dostarcza narzędzi do analityki biznesowej i zarządzania kadrami,
    \item centralizuje techniczne zarządzanie infrastrukturą systemową.
\end{itemize}

\section{Charakterystyka grup użytkowników}
W systemie zdefiniowano siedem dedykowanych ról, z których każda operuje na interfejsie dostosowanym do specyfiki wykonywanych zadań oraz urządzenia końcowego.

\subsection{Gość hotelowy (Interfejs tabletowy)}
Moduł przeznaczony dla osób zameldowanych w obiekcie, zlokalizowany na dedykowanych tabletach w pokojach hotelowych. Interfejs skupia się na zapewnieniu intuicyjnego dostępu do usług hotelowych, takich jak zamówienia gastronomiczne, zgłaszanie potrzeb serwisowych oraz bezpośrednia komunikacja z obsługą obiektu.

\subsection{ Obsługa hotelowa / Housekeeping (Interfejs mobilny)}
Narzędzie mobilne dedykowane personelowi sprzątającemu. System pozwala na bieżące monitorowanie planu pracy, otrzymywanie powiadomień o potrzebach gości oraz aktualizację statusów czystości pokoi w czasie rzeczywistym.

\subsection{Konserwator / Maintenance (Interfejs mobilny)}
Aplikacja mobilna przeznaczona dla działu technicznego. Moduł służy do zarządzania zgłoszeniami dotyczącymi usterek i awarii, umożliwiając szybką reakcję na problemy techniczne zgłaszane przez gości lub inny personel.

\subsection{Recepcja (Interfejs desktopowy)}
Stanowisko operacyjne służące do nadzoru nad bieżącym obłożeniem hotelu. Interfejs ten dostarcza kluczowych informacji o stanie pokoi, co pozwala pracownikom recepcji na sprawną realizację procesów zameldowania i wymeldowania.

\subsection{Kuchnia / Room service (Interfejs tabletowy)}
Panel operacyjny dedykowany pracownikom gastronomii. Interfejs koncentruje się na czytelnej prezentacji napływających zamówień posiłków, zarządzaniu ich kolejnością oraz informowaniu o zmianach statusu przygotowania dań.

\subsection{Menadżer (Interfejs desktopowy)}
Nadrzędny moduł zarządczy odpowiedzialny za optymalizację pracy całego obiektu. Menadżer posiada dostęp do danych kadrowych, finansowych oraz harmonogramów pracy, a także zajmuje się delegowaniem zadań i analityką ich wydajności.

\subsection{Administrator (Interfejs desktopowy)}
Moduł o najwyższym poziomie uprawnień technicznych, przeznaczony do zarządzania warstwą systemową. Administrator odpowiada za zarządzanie kontami użytkowników i definiowanie zakresu ich uprawnień w systemie French Riviera.